\chapter{Stroke}

\section*{Definition}
Sudden interruption of blood flow to the brain (ischemic stroke) or bleeding within or around the brain (hemorrhagic stroke), causing rapidly developing neurological deficits. Both require emergency assessment.

\section*{What Happens in Your Body}
Ischemic stroke results from an arterial clot or critical narrowing; hemorrhagic stroke results from rupture of a blood vessel. Brain injury progresses with time, so symptoms require immediate assessment even if they improve.

\section*{Causes}
\begin{itemize}
  \item Ischemic causes: atrial fibrillation or other cardiac emboli, carotid or intracranial atherosclerosis, and small-vessel disease
  \item Hemorrhagic causes: uncontrolled hypertension, cerebral aneurysm, arteriovenous malformation, anticoagulant-related bleeding, or head trauma
  \item Less common causes: arterial dissection, cerebral venous thrombosis, vasculitis, and inherited or acquired clotting disorders
\end{itemize}

\section*{When to See a Doctor}
Seek care immediately:
\begin{itemize}
  \item Sudden onset of any neurological deficit
  \item Time is brain --- call local emergency services immediately (Face drooping, Arm weakness, Speech difficulty, Time to call)
  \item Symptoms that resolve may still represent a transient ischemic attack and require urgent assessment
\end{itemize}

\section*{Related Symptoms}
\begin{itemize}
  \item Facial drooping or weakness
  \item Weakness or numbness of an arm or leg, especially on one side
  \item Difficulty speaking or understanding speech
  \item Sudden vision loss, double vision, dizziness, or trouble walking
  \item Sudden severe headache, especially with vomiting or altered alertness
\end{itemize}
\textbf{Important:} Stroke is a medical emergency requiring immediate evaluation.


\section*{Self-Care / Home Management}
\begin{itemize}
  \item Learn the FAST signs: Face drooping, Arm weakness, Speech difficulty, Time to call emergency.
  \item Every minute counts — do not wait to see if symptoms go away. Call emergency services immediately.
  \item Keep a list of your medications and medical conditions accessible for emergency responders.
  \item After a stroke, follow your rehabilitation plan closely for the best recovery.
  \item Control risk factors: blood pressure, diabetes, cholesterol, and stop smoking.
  \item Take prescribed blood thinners or other medications exactly as directed.
\end{itemize}

\section*{How Your Doctor Will Evaluate You}
\begin{itemize}
  \item Detailed neurological history includes onset pattern, progression, triggers, associated symptoms, and functional impact.
  \item Comprehensive neurological examination assesses mental status, cranial nerves, motor/sensory function, reflexes, coordination, and gait.
  \item Neuroimaging (CT, MRI) is often indicated to evaluate for structural, vascular, or demyelinating causes.
  \item Specialized testing (EEG for seizures, nerve conduction studies for neuropathy, lumbar puncture for meningitis) is guided by clinical suspicion.
\end{itemize}

\section*{Treatment Options}
\begin{itemize}
  \item Emergency brain imaging distinguishes ischemic from hemorrhagic stroke before clot-busting or antithrombotic treatment is considered.
  \item Eligible patients with acute ischemic stroke may receive intravenous thrombolysis within the approved time window; selected patients with large-vessel occlusion may benefit from mechanical thrombectomy, sometimes up to 24 hours after onset based on imaging.
  \item Hemorrhagic stroke requires specialist management, including blood-pressure control, reversal of anticoagulation when indicated, and neurosurgical or endovascular treatment for selected causes.
  \item Stroke-unit care, followed by physical, occupational, and speech therapy, supports recovery and secondary prevention.
\end{itemize}

\section*{References}
\begin{thebibliography}{99}
\bibitem{stroke:aha-stroke} Powers WJ, et al. ``AHA/ASA guideline for the early management of acute ischemic stroke.'' \emph{Stroke}. 2024;55(1):e1-e67.
\bibitem{stroke:uptodate-stroke} Caplan LR, et al. Clinical diagnosis of stroke. In: UpToDate, Post TW (Ed), UpToDate, Waltham, MA; 2025.
\bibitem{stroke:nejm-stroke} Campbell BCV, et al. ``Ischaemic Stroke.'' \emph{N Engl J Med}. 2023;389(16):1498-1510.
\bibitem{stroke:lancet-stroke} Feigin VL, et al. ``Global burden of stroke.'' \emph{Lancet}. 2024;403(10421):510-524.
\bibitem{stroke:bmj-stroke} Hankey GJ, et al. ``Stroke.'' \emph{BMJ}. 2024;386:e082543.
\bibitem{stroke:harrison-stroke} Jameson JL, et al. \emph{Harrison\'s Principles of Internal Medicine}. 21st ed. McGraw-Hill; 2022. Chapter 106: Ischemic Stroke.
\end{thebibliography}
